\section{What the agents actually built}
\label{sec:agentlogic}

A score on a frozen test set cannot distinguish two runs that arrive at the same
error by entirely different reasoning. One might choose its adaptive points by
ensemble disagreement and stop because validation error crossed the task's
threshold; another might relabel uniform sampling as ``adaptive'' and stop
because it ran out of rounds. Those are different research behaviours and the
same number.

Since that difference is much of what we wanted to measure, we measure it
directly.

\subsection{How it is measured}

Two extractors run over each finished workspace. Both are \textbf{LLM-free and
execution-free}: they parse the agent's Python with the standard \code{ast}
module and match patterns scoped to the role a file plays. The whole archive
therefore re-extracts identically, for nothing, as often as we like.

Scope is load-bearing. A whole-workspace search for \code{subprocess} is true of
nearly every run and tells you nothing, so each question is answered only from
the code that plays the relevant role --- the masking rule from
\code{predict.py} and the solver wrapper, the self-test from whatever actually
invokes \code{predict.py} as a subprocess --- and files that are clearly not
production code (tests, smoke checks, repro scripts) are excluded. Every answer
carries \code{file:line} evidence, so any cell of the resulting table can be
audited back to the line that produced it.

The vocabulary is anchored on the typed action space of the scripted levels, so
that a scripted policy, a typed picker and a from-scratch agent all land in the
same table. Anything unmatched stays \code{unknown} rather than being forced
into a category. These are measurements placed beside the contract, never
gates.

\begin{table}[t]
\caption{What each cell chose, and whether its code agrees with its prose.
\emph{Agr.} is the share of a cell's replicates on its modal chain of methods
--- method reproducibility, which is not the same thing as score
reproducibility. \emph{Switch.} counts replicates whose acquisition criterion
changed between rounds. \emph{M-inf.} counts replicates whose point-choosing
function actually calls the surrogate. The last column is the modal verdict of
comparing the criterion each run names against the one its code computes.}
\label{tab:logic}
\footnotesize
\setlength{\tabcolsep}{3.4pt}
\begin{tabular}{@{}lllcccl@{}}
\toprule
\textbf{Cell} & \textbf{Model} & \textbf{Acquisition} &
\textbf{Agr.} & \textbf{Switch.} & \textbf{M-inf.} & \textbf{Code vs prose} \\
\midrule
\Ltwo{}-\hn{cursor} & kernel ridge & \code{uncertainty\_gp} & 0.20 & 0/2 & --- & unverifiable 5/5 \\
\Ltwo{}-\hn{dspy}   & kernel ridge & \code{residual\_ucb}   & 0.10 & 1/9 & 10/10 & partial 8, agree 2 \\
\Ltwo{}-\hn{pi}     & kernel ridge & \code{residual\_ucb}   & 0.20 & 1/5 & --- & unverifiable 5/5 \\
\Ltwo{}-\hn{ursa}   & ridge / k.r. & \code{uncertainty}     & 0.38 & 0/3 & 2/5 & partial 3, \textbf{mismatch 2} \\
\midrule
\Lthree{}-\hn{cursor} & MLP (Torch) & \code{uncertainty\_ensemble} & 0.20 & 0/5 & 5/5 & \textbf{agree 4}, partial 1 \\
\Lthree{}-\hn{dspy}   & MLP (Torch) & \code{uncertainty\_ensemble} & 0.10 & 0/8 & 8/10 & partial 9, mismatch 1 \\
\Lthree{}-\hn{pi}     & MLP (Torch) & \code{uncertainty\_ensemble} & 0.20 & 0/3 & 5/5 & \textbf{agree 3}, partial 2 \\
\Lthree{}-\hn{ursa}   & MLP (Torch) & \code{uncertainty}           & 0.17 & 0/2 & 4/5 & partial 4, mismatch 1 \\
\bottomrule
\end{tabular}
\end{table}

\subsection{Method choice tracks the access level, not the substrate}

Table~\ref{tab:logic} shows a split that no outcome metric reveals.

At \Ltwo{}, the primary model is \textbf{kernel ridge in 19 of the 25
replicates} with a classifiable primary --- modal in three of the four cells,
and tied first in the fourth. Kernel ridge is also what the library's own model
zoo selects in all six runs of the \Lzero{} arm. At \Lthree{} the primary model
is \textbf{a hand-written PyTorch MLP in all 25}, without a single exception.

Acquisition splits the same way. \Ltwo{} agents gravitate to residual-UCB and
GP variance --- the two strategies the \Ltwo{} prompt names as available in the
library --- while \Lthree{} agents converge on ensemble disagreement, which no
prompt mentions anywhere.

So: four independent agent frameworks, given a library, settle on its defaults
or tie with them. Given nothing, all four rediscover a different method family,
and two of them beat the library with it (Section~\ref{sec:vs-baseline}). That
is the mechanism behind the level effect of Section~\ref{sec:l2l3}, and it is
why we read that effect as anchoring rather than as the library being poor
code.

\subsection{Method reproducibility is poor even where score reproducibility is good}

Chain agreement --- the share of a cell's replicates that chose the same chain
of methods --- never exceeds $0.38$, and is $0.10$ in both ten-replicate cells.
The same agent, given the same prompt and the same model, produces a different
pipeline nearly every time.

\Lthree{}-\hn{cursor} is the sharpest case: five distinct method chains in five
replicates, while its error stayed inside $[0.062, 0.212]$ and every artifact
was physically valid. A benchmark that replicates a cell and reports only the
spread of scores would conclude the agent is consistent. It is consistently
\emph{good}. It is not doing the same thing twice, and if what you care about is
whether an agent has a reliable method rather than a reliable outcome, those are
very different findings.

\subsection{Adaptivity is nominal: one fixed rule, executed $n$ times}

The task mandates up to three adaptive rounds whose points are chosen from the
current model or data. In \textbf{2 of the 37} runs where the criterion could be
classified across rounds did that criterion actually change between rounds. The
rest ran one rule three times.

The crude failure --- a stated criterion that is simply randomness --- occurs in
no cell at all. It is the subtler one, non-reactive adaptivity, that is
near-universal. That is consistent with the control arm of
Section~\ref{sec:control-arm} finding little for a reactive criterion to earn at
this budget: the agents may be leaving nothing much on the table.

\subsection{The code often does not match the prose, and a library hides it}

Comparing the criterion each run \emph{states} against the one its code
\emph{computes}, at the level of the criterion family, gives: \code{partial} 27,
\code{unverifiable\_from\_code} 14, \code{agree} 9, \code{mismatch} 4, of the 54
runs with a verdict. \code{agree} is the modal verdict in only two cells, both
\Lthree{}. The four outright mismatches --- stated and implemented criteria with
nothing in common --- fall in the two substrates that also produce every invalid
artifact.

Seven \Ltwo{} runs name a method in their README or report for which the code
carries no evidence at all, against \textbf{zero such claims in any \Lthree{}
cell}.

More awkward for anyone trying to evaluate an agent: \Ltwo{}-\hn{cursor} and
\Ltwo{}-\hn{pi} are \code{unverifiable\_from\_code} in all ten of their runs,
and whether their point selection is model-informed is simply unknown. Selection
happens inside the library, so no function in the workspace chooses points at
all. This is a structural consequence of giving an agent scaffolding rather than
evasion by the agent --- but it means that \textbf{the more scaffolding an
evaluation provides, the less of the agent's reasoning remains inspectable}.
Where a chooser can be classified at all it usually does call the surrogate: 22
of 25 such runs at \Lthree{} and 12 of 15 at \Ltwo{}.

Relatedly, \hn{ursa}'s rounds largely did not record validation error against
baseline at the moment of choosing (median evidence-grounding $0.33$ at \Ltwo{}
and $0.17$ at \Lthree{}, against $1.0$ in the six other cells), so its stopping
decisions cannot be falsified from its own artifacts.

\subsection{Where the rest of the solution agrees, and where it only nearly does}

Twelve further dimensions ask how each run solved the remainder of the task ---
who owns the solver's single global environment, what insulates one solve from
the next, how the triangular mesh reaches the fixed output grid, and so on. The
full cross-comparison is in Appendix~\ref{app:shape}.

One prohibition held outright: \textbf{no \Lthree{} run hand-built a mesh},
26 of 26, which is the explicit anti-Delaunay constraint in the task text doing
its job.

Three dimensions look unanimous at the level of each cell's modal value, and a
modal value hides the minority. The per-run counts, over the 51 runs whose
workspace can be read, are: 50 enforce the pilot gate, 46 validate stored data
before counting a solve --- \emph{three claim to and provide no evidence} ---
and 42 of the 43 with a readable answer keep the test set out of every
model-fitting function. \emph{One does not.}

The mandatory constraint most widely sidestepped is the deliverable self-test.
The task requires re-running the documented \code{predict.py} command in a fresh
process before declaring done. Twenty-four runs satisfy that by \emph{importing}
the module instead, against 22 that actually run it as a subprocess. Code that
has not been re-run the way it is documented to be run is the precondition for
most of the interface failures this benchmark has ever recorded --- and it is
measurable only because the question is asked of the code that plays that role,
rather than of the workspace as a whole.
